\section{Problem Formulation}

\subsection{Sensor scheduling formulation}

Let $x_t \in \mathbb{R}^d$ denote the simulated physical state at time step
$t$. The state contains meteorological and blowing-snow variables, including
wind speed and direction, air temperature, relative humidity, air pressure,
solar radiation, snow surface temperature, snow particle statistics, and snow mass flux. Let
$a_t \in \{0,1\}^m$ denote the activation vector applied to the $m$ logical
channels at time step $t$. Each logical sensing channel is a simulation-level
unit with associated observations, power demand, startup demand, and warm-up
behavior.

\begin{table}[htbp]
  \centering
  \caption{Notation used for scheduler inputs, actions, and constraints.}
  \label{tab:notation_summary}
  \footnotesize
  \setlength{\tabcolsep}{4pt}
  \begin{tabular}{@{}p{0.18\linewidth}p{0.72\linewidth}@{}}
    \toprule
    Symbol & Meaning \\
    \midrule
    $x_t$ & Simulated physical state at time step $t$ \\
    $o_t$ & Scheduler observation built from observation histories, channel modes, channel sampling ages, the previous executed activation vector, budget/time features, and simulated warning signals \\
    $\mathcal{A}$ & Candidate action set satisfying the core-channel and optional-channel capacity rules \\
    $K$ & Number of candidate activation vectors in $\mathcal{A}$ \\
    $n_{\mathrm{opt}}$ & Number of optional channels, $|\mathcal{S}|$ \\
    $\mathcal{F}_t$ & Feasible activation-vector set satisfying the steady-state and startup-peak power limits at time step $t$ \\
    $r_t$ & Remaining required activation time for the previous executed activation vector \\
    $a^{(k)}$ & Candidate activation vector indexed by $k$ \\
    $\mathcal{I}_t$ & Feasible action indices $\{k:a^{(k)}\in\mathcal{F}_t\}$ \\
    $u_t$ & Proposed action index sampled or selected from $\mathcal{I}_t$ \\
    $\bar a_t=a^{(u_t)}$ & Proposed activation vector \\
    $a_t$ & Executed activation vector after the minimum dwell-time constraint \\
    \bottomrule
  \end{tabular}
\end{table}

\begin{table}[htbp]
  \centering
  \caption{Notation used for forecast losses, supervision, and policy training.}
  \label{tab:notation_objectives}
  \scriptsize
  \setlength{\tabcolsep}{4pt}
  \begin{tabular}{@{}p{0.18\linewidth}p{0.72\linewidth}@{}}
    \toprule
    Symbol & Meaning \\
    \midrule
    $L_{\mathrm{step}}$ & Per-step multi-step forecast loss with component-wise clipping at 10 \\
    $L_{\mathrm{train}}$ & Condition-weighted, calibration/validation-normalized training loss used in the policy reward \\
    $L_c$ & Event-type-specific forecast loss for event type $c$ \\
    $s_c$ & Median calibration/validation event-type-specific forecast loss across candidate fixed channel subsets for $c\in\{\mathrm{particle},\mathrm{flux},\mathrm{thermal}\}$ \\
    $s_{\mathrm{calm}}$ & Median overall calibration/validation forecast loss across the same candidate fixed channel subsets; default normalization factor for calm time steps \\
    $L_{\mathrm{macro}}$ & Macro-averaged normalized forecast loss with equal weights for the three event types \\
    $L_{\rho}$ & Weighted normalized forecast loss used in Proposition~\ref{prop:specialist_bottleneck} \\
    $\mathcal{V}_{\mathrm{AoI}}$ & Observed variables whose ages enter the theoretical AoI objective; not necessarily the forecast-target set $\mathcal{Y}$ \\
    $\pi_{\text{label}}$ & Theoretical privileged policy with access to condition labels \\
    $\pi_s$ & Fixed-channel policy that selects optional channel $s$ \\
    $c_t$ & Current condition label used for training-loss weighting on the policy-training partition and offline loss computation and grouping on the test partition; never a policy or critic input at decision time \\
    $c_t^{\mathrm{sup}}$ & Future-condition supervision label: the first non-calm condition label in $[t,t+16]$, or calm if none occurs; used only for behavior cloning and auxiliary classification during policy training \\
    $\ell_t$ & Indicator of a valid supervised target action; $\ell_t=1$ at every PPO step in the reported stride-one configuration \\
    $\eta_t$ & Indicator of an available future-condition supervision label for the auxiliary classifier; $\eta_t=\ell_t=1$ at every PPO step in the reported stride-one configuration \\
    $S_t$ & Normalized Hamming distance $m^{-1}\lVert a_t-a_{t-1}\rVert_1$; fraction of executed channel activation bits that change \\
    $W_t$ & Number of warm-up interruptions at time step $t$ \\
    \bottomrule
  \end{tabular}
\end{table}


The sensing system contains a mandatory meteorological core channel and a set of
optional channels. The core channel remains active because it supplies the
minimum measurements needed by the downstream forecast task. At most one
optional channel can be active at each time step. The candidate action set is
\begin{equation}
  \mathcal{A}
  =
  \left\{a\in\{0,1\}^{m}\;\middle|\;
    a_{\mathrm{met}}=1,\;
    \sum_{i\in\mathcal{S}}a_i\leq1
  \right\}.
  \label{eq:candidate_action_set}
\end{equation}

Given the previously executed activation vector $a_{t-1}$, the feasible set
contains the candidate actions that satisfy the steady-state and startup-peak power
limits.
\begin{equation}
  \begin{aligned}
  \mathcal{F}_t
  =\bigg\{a\in\mathcal{A}\;\bigg|\;&
    \sum_i p_i a_i\leq B,\\[-2pt]
  & \sum_{i:a_{t-1,i}=1} p_i a_i
    +\sum_{i:a_{t-1,i}=0}
      \max\!\left(p_i,p_i^{\mathrm{peak}}\right)a_i
    \leq B^{\mathrm{peak}}
  \bigg\}.
  \label{eq:feasible_action_set}
  \end{aligned}
\end{equation}

Here, $\mathcal{S}$ denotes the optional-channel index set. The steady-state powers
$p_i$ and budget $B$ use a common normalized unit. The startup-peak powers
$p_i^{\mathrm{peak}}$ and budget $B^{\mathrm{peak}}$ use a second common
normalized unit. An active channel contributes its steady-state power to the
startup-power check. A newly activated channel contributes the larger of its
steady-state and startup-peak power. Together with \Cref{eq:candidate_action_set}, these constraints permit a fixed
schedule to select one optional channel while preventing more than one optional
channel from operating at the same time.

The finite candidate set is indexed as
$\{a^{(1)},\ldots,a^{(K)}\}=\mathcal{A}$. At time step $t$, the feasible
indices are
\[
  \mathcal{I}_t=\{k:a^{(k)}\in\mathcal{F}_t\}.
\]
During training, the policy samples an index from $\mathcal{I}_t$. During
evaluation, it selects the highest-scoring feasible index. Let $u_t$ denote the
resulting index and let $\bar a_t=a^{(u_t)}$ be the proposed activation vector.
Let $r_t$ denote the remaining required activation time for the previous
executed activation vector. The executed activation vector is
\begin{equation}
  a_t=
  \begin{cases}
    a_{t-1}, & r_t>0,\\
    \bar a_t, & r_t=0.
  \end{cases}
  \label{eq:minimum_activation_rule}
\end{equation}

The feasible-action mask restricts the policy distribution to
$\mathcal{I}_t$ and enforces the steady-state and startup-peak power limits. The
execution rule retains $a_{t-1}$ until its minimum dwell time has been
completed. Each executed activation vector therefore remains active for at least
$D_{\min}$ time steps.

\subsection{Conditions under which optional-channel adaptation reduces forecast loss}

The sensing configuration with a mandatory meteorological core channel and at most one
active optional channel is a special case of a broader capacity-constrained
scheduling problem. Let $n_{\mathrm{opt}}=|\mathcal{S}|$ denote the number of
optional channels, of which the scheduler may select at most
$q<n_{\mathrm{opt}}$ at each time step. For an operating
condition $c\in\mathcal{C}$, let
$S_c^\star\subseteq\mathcal{S}$ denote an optimal optional-channel subset, with
$|S_c^\star|\leq q$.

Let $\widetilde L_c(\pi)$ denote the normalized forecast loss of policy $\pi$
under condition $c$. The weighted normalized forecast loss is
\[
L_{\rho}(\pi)
=
\sum_{c\in\mathcal{C}}
\rho_c\widetilde L_c(\pi),
\qquad
\rho_c>0,
\qquad
\sum_{c\in\mathcal{C}}\rho_c=1.
\]
The empirical macro-averaged normalized forecast loss in
\Cref{eq:macro_loss} uses
$\widetilde L_c(\pi)=L_c(\pi)/s_c$ and assigns equal weight to the particle,
flux, and thermal event types.

Dynamic scheduling can reduce forecast loss when no fixed optional-channel
subset is optimal for every condition with positive evaluation weight. The
advantage requires positive excess forecast loss when a fixed subset is
mismatched to the operating condition. Predictive information in the scheduler
input can then support channel changes before or during a condition transition.
The following proposition gives a sufficient condition for this advantage.

\begin{proposition}[Dynamic value with one active optional channel]
\label{prop:specialist_bottleneck}
Consider a scheduling problem that permits at most one active optional channel.
All policies are subject to the same startup and minimum dwell-time
constraints. For any $s\in\mathcal{S}$, let $\pi_s$ denote the fixed-channel
policy that selects $s$ under every operating condition. Suppose that at least
two positive-weight conditions have different optimal optional channels. The
corresponding condition intervals are long enough for the theoretical privileged
policy $\pi_{\text{label}}$, which has condition-label access, to complete
feasible channel changes under the common
execution constraints.

Transition losses are included in the normalized condition-specific losses.
For every fixed policy $\pi_s$ and every condition $c$,
\[
\widetilde L_c(\pi_s)
\geq
\widetilde L_c(\pi_{\text{label}}).
\]
For each positive-weight condition whose optimal channel differs from $s$,
\[
\widetilde L_c(\pi_s)
-
\widetilde L_c(\pi_{\text{label}})
\geq
\Delta_c
>
0.
\]
Then every fixed policy $\pi_s$ that always selects one optional channel
$s\in\mathcal{S}$ satisfies
\[
L_\rho(\pi_s)
>
L_\rho(\pi_{\text{label}}).
\]
\end{proposition}

A proof is provided in \ref{apx:specialist_bottleneck_proof}, and
\Cref{fig:specialist_bottleneck} illustrates the sufficient condition. The
proposition assumes that all policies follow the same startup and minimum
dwell-time constraints. It does not provide a performance bound for the
PD-PPO scheduler or any other implemented policy. The empirical privileged true-label baseline in
Supplementary Table S11 is not assumed to implement the theoretical
policy $\pi_{\text{label}}$ or to attain the minimum forecast
loss.

\begin{figure}[H]
  \centering
  \resizebox{\linewidth}{!}{\begin{tikzpicture}[
  x=1cm,
  y=1cm,
  >=Latex,
  font=\sffamily\fontsize{7.8}{9.2}\selectfont,
  panel/.style={rounded corners=2.2mm,line width=0.8pt,minimum width=4.05cm,
    minimum height=5.00cm,inner sep=0pt},
  title/.style={font=\sffamily\fontsize{8.8}{10.3}\selectfont\bfseries,
    text=pdcharcoal,align=left},
  badge/.style={font=\sffamily\fontsize{9.2}{10.8}\selectfont\bfseries,
    text=pdcharcoal,align=center},
  subtitle/.style={font=\sffamily\fontsize{7.4}{8.6}\selectfont,
    text=pdslate,align=center,text width=3.65cm},
  matrixcell/.style={rounded corners=0.8mm,draw=white,line width=0.5pt,
    minimum width=0.72cm,minimum height=0.46cm,align=center},
  rowlabel/.style={font=\sffamily\bfseries,text=pdslate,align=right},
  item/.style={rounded corners=1.2mm,draw=white,line width=0.5pt,
    minimum width=3.45cm,minimum height=0.53cm,align=center},
  flow/.style={-Latex,line width=1.15pt,draw=pdslate},
  note/.style={font=\sffamily\fontsize{7.2}{8.5}\selectfont,text=pdslate,
    align=center,text width=3.65cm}
]
  \definecolor{pdblue}{HTML}{0072B2}
  \definecolor{pdteal}{HTML}{009E73}
  \definecolor{pdorange}{HTML}{E69F00}
  \definecolor{pdslate}{HTML}{56606A}
  \definecolor{pdcharcoal}{HTML}{28323C}
  \definecolor{pdlight}{HTML}{EEF2F5}

  \node[panel,fill=pdblue!5,draw=pdblue!65] (panelA) at (2.05,0) {};
  \node[panel,fill=pdorange!6,draw=pdorange!72] (panelB) at (6.55,0) {};
  \node[panel,fill=pdteal!6,draw=pdteal!72] (panelC) at (11.05,0) {};

  \node[badge] at (0.30,2.10) {A};
  \node[title,anchor=west] at (0.60,2.10) {Condition optima};
  \node[subtitle] at (2.05,1.45)
    {$n_{\mathrm{opt}}=3$, $q=1$; best channel\\depends on the condition};
  \node[rowlabel] at (0.75,0.72) {$c\backslash s$};
  \node[rowlabel] at (0.75,0.18) {$c_1$};
  \node[rowlabel] at (0.75,-0.36) {$c_2$};
  \node[rowlabel] at (0.75,-0.90) {$c_3$};
  \node[rowlabel,align=center] at (1.42,0.72) {$s_1$};
  \node[rowlabel,align=center] at (2.22,0.72) {$s_2$};
  \node[rowlabel,align=center] at (3.02,0.72) {$s_3$};
  \node[matrixcell,fill=pdteal!88,text=white,font=\sffamily\bfseries]
    at (1.42,0.18) {$s_{c_1}^{\star}$};
  \node[matrixcell,fill=pdlight,text=pdslate] at (2.22,0.18) {$\times$};
  \node[matrixcell,fill=pdlight,text=pdslate] at (3.02,0.18) {$\times$};
  \node[matrixcell,fill=pdlight,text=pdslate] at (1.42,-0.36) {$\times$};
  \node[matrixcell,fill=pdteal!88,text=white,font=\sffamily\bfseries]
    at (2.22,-0.36) {$s_{c_2}^{\star}$};
  \node[matrixcell,fill=pdlight,text=pdslate] at (3.02,-0.36) {$\times$};
  \node[matrixcell,fill=pdlight,text=pdslate] at (1.42,-0.90) {$\times$};
  \node[matrixcell,fill=pdlight,text=pdslate] at (2.22,-0.90) {$\times$};
  \node[matrixcell,fill=pdteal!88,text=white,font=\sffamily\bfseries]
    at (3.02,-0.90) {$s_{c_3}^{\star}$};
  \node[note] at (2.05,-1.85)
    {No single optional channel\\is best for all conditions};

  \node[badge] at (4.80,2.10) {B};
  \node[title,anchor=west] at (5.10,2.10) {Fixed allocation};
  \node[subtitle] at (6.55,1.45) {Example: $s=s_2$\\for every condition};
  \node[item,fill=pdorange!23,text=pdcharcoal] at (6.55,0.60)
    {$c_1\rightarrow s_2$\quad excess $+\Delta_1$};
  \node[item,fill=pdteal!16,text=pdcharcoal] at (6.55,-0.02)
    {$c_2\rightarrow s_2$\quad condition best};
  \node[item,fill=pdorange!23,text=pdcharcoal] at (6.55,-0.64)
    {$c_3\rightarrow s_2$\quad excess $+\Delta_3$};
  \node[note] at (6.55,-1.85)
    {Every fixed $s$\\has at least one mismatch};

  \node[badge] at (9.30,2.10) {C};
  \node[title,anchor=west] at (9.60,2.10) {Label-aware policy};
  \node[subtitle] at (11.05,1.35) {Condition-label access\\guides channel changes};
  \node[item,fill=pdteal!16,text=pdcharcoal] at (11.05,0.60)
    {$c_1\rightarrow s_1=s_{c_1}^{\star}$};
  \node[item,fill=pdteal!16,text=pdcharcoal] at (11.05,-0.02)
    {$c_2\rightarrow s_2=s_{c_2}^{\star}$};
  \node[item,fill=pdteal!16,text=pdcharcoal] at (11.05,-0.64)
    {$c_3\rightarrow s_3=s_{c_3}^{\star}$};
  \node[note] at (11.05,-1.85)
    {Transition losses enter\\each condition-specific loss};

  \draw[flow] (panelA.east) -- (panelB.west);
  \draw[flow] (panelB.east) -- (panelC.west);

  \node[rounded corners=1.6mm,draw=pdslate!55,fill=white,line width=0.7pt,
    minimum width=13.05cm,minimum height=0.78cm,align=center,
    font=\sffamily\fontsize{9.0}{10.6}\selectfont,text=pdcharcoal]
    at (6.55,-3.25)
    {$\displaystyle
      L_{\rho}(\pi_s)-
      L_{\rho}(\pi_{\text{label}})
      \geq \sum_{c\in\mathcal{C}_{\mathrm{mis}}(s)}
      \rho_c\Delta_c > 0$};
\end{tikzpicture}
}
  \caption{Sufficient condition for the value of adaptive scheduling with one
  active optional channel. Different operating conditions favor different
  optional channels, which causes every fixed policy that always selects one
  optional channel to incur excess
  loss for at least one condition. The theoretical privileged policy can change channels
  when the startup and minimum dwell-time constraints allow a feasible
  transition. The lower panel gives the resulting bound on weighted normalized
  forecast loss. The set $\mathcal{C}_{\mathrm{mis}}(s)$ contains the
  positive-weight conditions whose optimal channel differs from fixed channel
  $s$. The quantity $\Delta_c$ is the corresponding normalized excess loss.}
  \label{fig:specialist_bottleneck}
\end{figure}

\paragraph{Relation between the simulation and the proposition}
The simulation was designed to instantiate condition-dependent channel optima
of the type considered in \Cref{prop:specialist_bottleneck}. During particle,
flux, and thermal intervals,
the laser disdrometer channel, FC4 snow mass-flux channel, and infrared
snow-surface-temperature channel, respectively, provide the most reliable measurements of the
target variables emphasized for the corresponding event type. The event-type
decomposition on the test partition in
\Cref{tab:clean_regime_decomposition} evaluates whether the PD-PPO scheduler
reduces loss under these event-specific measurement requirements and improves
the macro-averaged normalized forecast loss. It does not verify the proposition's
inequality for every $\pi_s$ and the theoretical $\pi_{\mathrm{label}}$ in every
condition. The proposition does not require every event type to have lower loss
in every seed.

\subsection{Forecast loss and macro-averaged evaluation across event types}

The frozen primary forecaster receives a window of sample-and-hold values and
observation masks and predicts selected target variables at horizons
$k=1,\ldots,H$. Let $\hat{y}_{t+k,j}$ and $y_{t+k,j}$ denote the prediction
and target, respectively, for variable $j\in\mathcal{Y}$ at horizon $k$. The
condition label at time step $t$ is
\begin{equation}
  c_t\in\{\text{calm},\text{particle},\text{flux},\text{thermal}\}.
\end{equation}

Let $w_k$ denote the horizon weight, $\alpha_{j,c_t}$ the condition-specific
target weight, and $\sigma_j$ the scale of target variable $j$. The per-step
forecast loss is
\begin{equation}
  L_{\text{step}}(t)
  =
  \frac{
  \sum_{k=1}^{H}\sum_{j \in \mathcal{Y}}
  w_k \alpha_{j,c_t}
  \min\!\left\{10,\,
  \left| \hat{y}_{t+k,j} - y_{t+k,j} \right|/\sigma_j\right\}
  }{
  \sum_{k=1}^{H}\sum_{j \in \mathcal{Y}}
  w_k \alpha_{j,c_t}
  } .
  \label{eq:step_forecast_loss}
\end{equation}

The target scales satisfy $\sigma_j>0$, and all horizon and target weights are
nonnegative. The denominator is positive for every scoreable time step. Each
scale-normalized absolute error is capped at 10 before weighting and aggregation.

The experiment uses uniform horizon weights and the predeclared
condition-specific target weights in Supplementary Table S4. The condition
label determines which target weights are applied during policy training and
evaluation. It is not included in the scheduler input. Losses for PD-PPO and
all baselines are computed using the same target scales, target weights, and
condition labels.

An overall average can be influenced disproportionately by event frequency and
by differences in raw loss magnitude across event types. The macro-averaged normalized
forecast loss addresses these effects by normalizing each event-type-specific
loss with calibration/validation results from feasible fixed channel subsets and assigning
equal weight to the three event types. Let
$c\in\{\mathrm{particle},\mathrm{flux},\mathrm{thermal}\}$ index the event
types, and let $\mathcal{T}_c$ denote the evaluated test time steps labeled as
$c$. For a policy $\pi$, the event-type-specific forecast loss is
\begin{equation}
  L_c(\pi)
  =
  \frac{1}{|\mathcal{T}_c|}
  \sum_{t\in\mathcal{T}_c}
  L_{\text{step}}(t;\pi).
\end{equation}

For each event type $c$, the normalization factor $s_c$ is the median
event-type-specific forecast loss across the feasible fixed channel subsets
evaluated on the calibration/validation partition. The macro-averaged normalized forecast
loss is
\begin{equation}
  L_{\text{macro}}(\pi)
  =
  \frac{1}{3}
  \sum_{c\in\{\mathrm{particle},\mathrm{flux},\mathrm{thermal}\}}
  \frac{L_c(\pi)}{s_c}.
  \label{eq:macro_loss}
\end{equation}

This metric is reported when $|\mathcal{T}_c|>0$ and $s_c>0$ for all three
event types. Lower values indicate better forecasting performance. The metric
compares policies across equally weighted event types after accounting for the
median calibration/validation loss of feasible fixed channel subsets.

\subsection{Non-equivalence of forecast loss, AoI, and covariance objectives}

Forecast loss, AoI, and covariance trace evaluate different properties of a
sensor schedule. Let $\pi$ be a feasible scheduling policy and let
$\gamma\in(0,1]$ denote a discount factor. The same horizon $T$ and temporal
weights $\gamma^t$ are used for all three objectives. The discounted forecast objective is
\begin{equation}
  J_{\mathrm{fcst}}(\pi)
  =
  -\mathbb{E}_{\pi}
  \left[
    \sum_{t=0}^{T-1}\gamma^t
    L_{\mathrm{step}}(t;\pi)
  \right].
  \label{eq:forecast_objective}
\end{equation}

For a separate model-based estimator with posterior covariance $P_t$, let
$\mathcal{V}_{\mathrm{AoI}}$ denote the observed variables whose ages enter an
AoI objective. This set need not equal the forecast-target set $\mathcal{Y}$.
The AoI and covariance objectives are
\begin{equation}
  J_{\mathrm{AoI}}(\pi)
  =
  -\mathbb{E}_{\pi}\!\left[
    \sum_{t=0}^{T-1}\gamma^t
    \sum_{j\in\mathcal{V}_{\mathrm{AoI}}}\mathrm{age}_{t,j}
  \right],
  \qquad
  J_{\mathrm{cov}}(\pi)
  =
  -\mathbb{E}_{\pi}\!\left[
    \sum_{t=0}^{T-1}\gamma^t
    \mathrm{tr}(P_t)
  \right].
\end{equation}

The three objectives can induce different orderings over the same feasible
schedules. Such a difference can occur when sensor variables have unequal
downstream forecast value, even if one schedule produces fresher observations
and a lower estimator covariance.

\begin{proposition}[Non-equivalence of forecast loss, AoI, and covariance objectives]
\label{prop:non_equivalence}
Forecast loss, covariance trace, and per-variable AoI objectives can induce
different orderings over feasible sensor schedules. In particular, there exist
a common horizon $T$, a common discount factor $\gamma$, and two feasible
policies $\pi_1$ and $\pi_2$ such that
\begin{equation}
  J_{\mathrm{AoI}}(\pi_1) > J_{\mathrm{AoI}}(\pi_2),
  \qquad
  J_{\mathrm{cov}}(\pi_1) > J_{\mathrm{cov}}(\pi_2),
\end{equation}
and
\begin{equation}
  J_{\mathrm{fcst}}(\pi_2) > J_{\mathrm{fcst}}(\pi_1).
\end{equation}
Equivalently, $\pi_1$ has lower AoI and lower covariance trace, and
$\pi_2$ has strictly lower forecast loss.
\end{proposition}

The proof in \ref{apx:non_equivalence_proof} takes the common horizon $T=1$
and considers a single-target, single-horizon special case of
\Cref{eq:step_forecast_loss}. A policy that refreshes the variable with greater
influence on the forecast target has strictly lower forecast loss in the
constructed example, while its total AoI and covariance trace are both higher.
This ordering shows that forecast loss must be evaluated directly when sensor
variables have unequal downstream forecast value.
The AoI-priority, round-robin, and random baselines provide non-learned dynamic
comparisons. The matched scalar-objective PPO configurations replace the scalar
scheduling objective while retaining the shared supervision and architecture.

The proposition establishes that the objectives can rank feasible schedules
differently. It does not determine the performance ordering of policies learned
from finite training data. The comparisons in
\Cref{sec:matched_controls,sec:objective_learner_controls} examine whether the
forecast-loss, AoI, and bounded diagonal uncertainty objectives produce
measurable differences under matched PPO configurations. The empirical uncertainty proxy is not the
posterior covariance trace used in the proposition.
